%\baselineskip11.8pt plus.2pt minus.2pt


%\normalsize
%\vskip\baselineskip
\advance\baselineskip by-.1pt

\niebowspis{Niebo w~styczniu}

\niebowtyt{Niebo w~styczniu}

\vspace*{-4pt}{\mathcode`\,="013B
Od początku roku dzień wydłuża się zarówno o~świcie, jak i~o~zmierzchu. Jednak tempo wydłużania dnia zwiększa się wyraźnie dopiero około 20 stycznia, gdy Słońce przekroczy równoleżnik $-20^{\circ}$ deklinacji w~drodze na północ. Od~tego momentu na naszych szerokościach geograficznych dnia przybywa o~ponad 2~godziny na miesiąc. 

Styczeń zacznie się dobrą widocznością Księżyca i~planet na wieczornym niebie. Na początku roku ekliptyka przyjmuje korzystne nachylenie do widnokręgu po zachodzie Słońca i~zbliża się do horyzontu nad ranem. Srebrny Glob przeszedł przez nów jeszcze w~grudniu, wędrując przy tym pod ekliptyką, i~od drugiego dnia miesiąca jego sierp ozdobi niebo po zmierzchu, na którym po zachodniej stronie  widoczne są planety Wenus i~Saturn.

\looseness-1{\spaceskip.33em minus.05emDruga planeta od Słońca przez cały miesiąc przejdzie od gwiazdozbioru Koziorożca przez Wodnika do Ryb, osiągając 10 stycznia maksymalną elongację wschodnią, wynoszącą $47^{\circ}$. W~tym czasie jasność Wenus zwiększy się do $-4,5^m$, a~średnica jej tarczy przekroczy $32''$ przy fazie zmniejszającej się do $37\%$. Po~drodze, 18 stycznia, Wenus minie Saturna, do którego zbliży się na $2^{\circ}$. Planeta wyraźnie poprawi swoje warunki obserwacyjne, wznosząc się o~zmierzchu na ponad $30^{\circ}$. 3~stycznia $2^{\circ}$ na południe od Wenus przejdzie Księżyc w~fazie $16\%$.}

\vspace*{\parskip}
\begin{dwieszpalty}



Ostatnia z~widocznych gołym okiem planet zbliża się do marcowej koniunkcji ze Słońcem i~swojej równonocy. W~styczniu Saturn świeci blaskiem $+1,1^m$, prezentując tarczę o~średnicy $16''$. Jego pierścienie stają się coraz węższe i~coraz trudniej je dostrzec przez teleskop, szczególnie przy małych powiększeniach. 7~i~23~stycznia na tarczy planety pojawi się cień Tytana, jej największego księżyca, co da się dostrzec nawet przez niewielkie teleskopy. 4~stycznia Saturn zbliży się na mniej niż $30'$ do gwiazdy 7. wielkości 85 Aqr, która może przypominać jeden z~jego księżyców, ale warto pamiętać, że nawet Tytan świeci o~prawie $2^m$ słabiej. Tego samego dnia dojdzie do zakrycia obu ciał niebieskich przez Księżyc w~fazie $25\%$. Saturn zniknie za ciemnym brzegiem księżycowej tarczy około godziny 18:40 i~pojawi się po jej jasnej stronie godzinę później.

6~stycznia Księżyc przejdzie przez I~kwadrę, a~4 dni później dotrze do gwiazdozbioru Byka, zwiększając fazę do $80\%$. Noc z~9~na 10 stycznia Srebrny Glob zacznie jakieś $5^{\circ}$ od Plejad, by około godziny 3 zakryć ich południowo-wschodnią część z~gwiazdą Merope. Niestety do zakryć dojdzie nisko nad widnokręgiem. Następnej nocy Srebrny Glob minie Jowisza w~odległości $5^{\circ}$. Planeta powoli słabnie, do końca miesiąca zmniejszając blask do $-2,4^m,$ a~średnicę tarczy do $44''$.

13 stycznia przed północą naszego czasu naturalny satelita Ziemi osiągnie pełnię, zajmując pozycję niecałe $4^{\circ}$ od Polluksa, najjaśniejszej gwiazdy Bliźniąt, a~o~godzinie 5:30 zbliży się na $0,5^{\circ}$ do Marsa. Czerwona Planeta 12 stycznia przejdzie najbliżej Ziemi (96~mln~km), by 4~dni później znaleźć się w~opozycji do Słońca. 23 stycznia Mars zbliży się na $2^{\circ}$ do Polluksa.

Jak zawsze zimowe opozycje Marsa, mimo tego, że planeta wznosi się wysoko na niebie, należą do tych niekorzystnych, i~w~tym sezonie średnica jego tarczy nie przekroczy $15''$. Również jego jasność nie zrobi dużego wrażenia, osiągnie maksymalnie $-1,4^m$, czyli porównywalnie z~Syriuszem. Oczywiście przy tej jasności planeta nadal wyprzedzi pod tym względem wszystkie pozostałe naturalne ciała niebieskie w~swoim sąsiedztwie, poza Księżycem, ale zarówno jej średnica kątowa, jak i~jasność nie zachwycą tych, którzy pamiętają wygląd jej tarczy podczas tzw. Wielkich Opozycji w~2018 czy 2003~roku. Niestety dla nas wtedy planeta przechodzi przez opozycję latem i~wędruje nisko nad widnokręgiem. Jest to jedna z~tych rzeczy, których szczerze możemy zazdrościć mieszkańcom półkuli południowej naszej planety, mających Marsa wysoko na niebie podczas jego największych opozycji. 

Księżyc powędruje dalej i~16 stycznia wzejdzie $2^{\circ}$ od Regulusa w~Lwie, a~5~dni później rano, mając tarczę oświetloną w~połowie, pokaże się $0,5^{\circ}$ od Spiki w~Pannie. Styczniowy nów przypada 29. dnia miesiąca. Do tego czasu 25 stycznia oświetlona w~$20\%$ tarcza Księżyca zajmie pozycję około $3^{\circ}$ od Antaresa, najjaśniejszej gwiazdy Skorpiona. 

Jak co roku na początku stycznia promieniują meteory z~roju Kwadrantydów. Ich radiant położony jest na północ od Wolarza, a~zatem nigdy u~nas nie zachodzi. Kwadrantydy najlepiej obserwować rano, gdy ich radiant wznosi się na ponad $60^{\circ}$. Są to meteory o~średniej prędkości, zderzają się z~naszą atmosferą z~prędkością 41~km/s. W~okolicach maksimum aktywności można spodziewać się nawet 80 zjawisk na godzinę.

4~lutego przewidywane jest maksimum aktywności długookresowej gwiazdy zmiennej R~Leo. Może ona osiągnąć blask nawet $+4,4^m$, jest wtedy dostrzegalna gołym okiem. R~Leo znajduje się jakieś $5^{\circ}$ od Regulusa, niedaleko gwiazd 6. wielkości 18 i~19 Leonis, z~którymi można porównywać jej blask. Ma też wyraźnie widoczną wiśniową barwę. Łatwo ją zatem zidentyfikować przez lornetkę. Pod koniec stycznia gwiazda góruje około godziny 1 na wysokości $50^{\circ}$.

%\medskip
\rightline{\large
  \it  Ariel MAJCHER}
\end{dwieszpalty}

}\eject
\normalsize